\section{multi-turn CVR}
\label{sec:intro}

We will detail the construction of our \Mybenchmark dataset.

% \subsection{Overview of \Mybenchmark} 
% We introduce the Multi-turn Composed Video Retrieval Benchmark (\Mybenchmark), a novel benchmark especially designed to evaluate the history-aware compositional retrieval capabilities of CVR models. Unlike prior CVR benchmarks focused solely on single-turn modified instruction, \Mybenchmark  systematically simulates real-world interactive retrieval scenarios. Each session is structured as $(\textit{source video}, \textit{user feedback}_1, \textit{result}_1, \textit{user feedback}_2, \ldots, \textit{result}_{N-1}, \textit{user feedback}_N, \textit{target video})$, 
% where the user feedback progressively describes how the source video or intermediate results should be modified to reach the target video.

% We curate over 120,000 videos from the WebVid2M dataset and partition them into 600 clusters. Within each cluster, we sample $(\textit{source video}, \textit{target video})$ pairs as the initial pairs for each session, resulting in 100,000 initial video pairs. Through $N$-turn interaction generation and careful filtering, we finally obtain 4,000 gold-standard $N$-turn test sessions. To better simulate real-world user interactions, we further organize \Mybenchmark into five representative multi-turn dialogue scenarios:
% \begin{itemize}
%     \item \textbf{Progression}: The user continuously refines the previous result, progressively moving closer to the target video.
%     \item \textbf{Rollback}: The user rejects the most recent result, returns to an earlier turn, and introduces a new modification based on that state.
%     \item \textbf{Combination}: The user integrates attributes from the current reference and an earlier turn into a single new modification instruction.
%     \item \textbf{Correction}: The user corrects or further refines the previous result.
%     \item \textbf{Preference Change}: The user shifts their original preference, disregards the prior interaction history, and instead specifies a new target video.
% \end{itemize}

% Unlike existing datasets, \Mybenchmark places particular emphasis on \textbf{history-aware queries} (e.g., Progression, Rollback, and Combination), which make up the majority of tasks. This distribution better reflects the complexity of real-world retrieval scenarios, where users rarely express their full retrieval intent in a single step and instead progressively refine their requests through multi-turn interactions with the system. Meanwhile, \textbf{Correction} and \textbf{Preference Change} also occur naturally in realistic user interactions. In such cases, historical context may become distracting rather than helpful. This requires retrieval systems not only to understand multimodal compositional queries, but also to accurately infer the user’s current intent. It is worth noting that the average feedback length in \Mybenchmark~(xx.x words) is higher than that of existing CVR benchmarks. This is a deliberate design choice. Traditional CVR queries are often derived by directly comparing source and target captions, which tends to produce short and potentially ambiguous instructions, such as \textit{``Make it huge.''} In contrast, \Mybenchmark generates more detailed and natural modification instructions by simulating realistic user feedback, such as \textit{``Actually, I want a big firework show.''}

% To further illustrate the semantic breadth of the benchmark, Figure~\ref{fig:n} presents the hierarchical distribution of content categories in the multi-turn dialogues. We categorize the data along two dimensions. The first is the \textbf{content dimension}. \Mybenchmark encompasses four primary domains: \textit{Natural and Geographic Scenes}, \textit{Human-Centric World}, \textit{Object-Centric World}, and \textit{Abstract \& Creative Content}. These domains are further organized into fine-grained subcategories—ranging from \textit{Flora \& Fauna} to \textit{Visual Arts}—to ensure broad coverage of diverse visual scenes. The second is the \textbf{modification intent}. Along this dimension, \Mybenchmark is carefully balanced across three high-level classes: \textit{Scene-Level Modifications}, \textit{Entity-Level Modifications}, and \textit{Attribute \& Style Modifications}, which are further divided into 15 distinct subcategories, such as \textit{Action \& Event Modification}, \textit{Subject \& Object Modification}, and \textit{Camera \& Viewpoint Control}. This fine-grained categorization across both dimensions highlights the semantic richness of \Mybenchmark and reflects the complexity of user modification instructions in practical interactive retrieval scenarios.

\subsection{Design Principles} 
Multi-turn composed video retrieval differs fundamentally from conventional single-turn CVR. In real interactive retrieval scenarios, users rarely express their full retrieval intent in a single step; instead, they progressively refine, revise, or even shift their requests through multi-turn feedback. As a result, the target retrieval intent is not determined solely by the latest instruction, but by how the current feedback interacts with the preceding dialogue history. Therefore, our benchmark for multi-turn CVR evaluate not only compositional understanding of the current query, but also a model’s ability to interpret evolving user intent under history-dependent interactions.

Drawing inspiration from common interaction patterns in real-world search sessions, we identify five representative scenarios and organize the sessions in \Mybenchmark accordingly. These scenarios capture different ways in which user intent evolves over multi-turn interactions, ranging from progressive refinement to history revision and preference shift: 

\begin{itemize}
    \item \textbf{Progression}: The user incrementally refines the current retrieval result, with each new feedback preserving the core intent of previous turns and moving the search progressively closer to the target video. 
    \item \textbf{Rollback}: The user rejects the most recent retrieval result, returns to an earlier interaction state, and continues the search trajectory from that historical point with a new modification request. 
    \item \textbf{Combination}: The user composes a new request by combining attributes or constraints from the current reference with those from an earlier turn, requiring the system to jointly preserve information across non-adjacent turns. 
    \item \textbf{Correction}: The user locally adjusts or rectifies the latest retrieval result, typically by correcting an undesired attribute or refining a previously underspecified modification. 
    \item \textbf{Preference Change}: The user shifts to a new retrieval preference during the interaction, making part or all of the earlier dialogue history no longer relevant to the current target. 
\end{itemize}

In addition to covering diverse interaction scenarios, \Mybenchmark is designed to span a broad range of visual content and modification types. Otherwise, the evaluation may be dominated by a narrow set of homogeneous scenes or edit patterns, limiting its ability to reflect the diversity of practical retrieval requests. To this end, \Mybenchmark is organized along two complementary dimensions: a \textbf{content dimension}, which spans four primary domains---\textit{Natural and Geographic Scenes}, \textit{Human-Centric World}, \textit{Object-Centric World}, and \textit{Abstract \& Creative Content}---and a \textbf{modification-type dimension}, which covers three high-level classes---\textit{Scene-Level Modifications}, \textit{Entity-Level Modifications}, and \textit{Attribute \& Style Modifications}. This design enables history-aware retrieval to be evaluated under diverse visual and compositional conditions.



% \subsection{Data Generation Pipeline} 
% Our data generation pipeline is designed to construct multi-turn retrieval sequences that are realistic, sufficiently challenging, and diverse in interaction patterns. To this end, we generate the benchmark in three stages:  Initial Multi-turn Interaction Generation, Interaction Scenario Construction, and Data Cleaning and Review.

% \paragraph{Stage 1: Initial Multi-turn Interaction Generation.} 
% We begin with WebVid-2M and partition the videos into 600 semantic clusters based on caption embeddings. This restricts reference–target sampling to locally coherent semantic neighborhoods, making the resulting edits more natural and compositionally meaningful. For each reference video, we sample a target from the same cluster under a similarity-range constraint, excluding candidates that are either too similar or too dissimilar. This encourages pairs that are suitable for progressive multi-turn refinement rather than trivial matching or abrupt semantic jumps. Given each reference–target pair, we prompt an MLLM to generate a user-style editing instruction describing one or two salient modifications, leaving room for subsequent turns. To control difficulty, we then perform retrieval within the same cluster using a pretrained \task model. If the target can already be retrieved in a single turn, the pair is discarded as overly easy. Otherwise, the top-1 retrieved result is used as the next-turn reference, and the process is repeated by conditioning the MLLM on the accumulated dialogue history and the fixed target. This yields seed multi-turn interactions that progressively approach the target video.

% \paragraph{Stage 2: Interaction Scenario Construction.} 
% While the interactions constructed in Stage 1 naturally capture progressive refinement, real-world interactive retrieval often involves more complex behaviors beyond turn-by-turn accumulation. These scenarios cover not only continuous refinement toward a fixed target, but also history-revising behaviors such as reverting to an earlier turn, combining constraints from multiple turns, correcting previous instructions, or shifting to a new target altogether. Based on the initial multi-turn interactions, we leverage GPT-4o to rewrite the dialogues according to these predefined scenarios. Specifically, rollback rewrites the current instruction by referring back to a selected historical turn, combination integrates attributes from both the current and an earlier turn, correction introduces self-revising language to refine or rectify the latest result, and preference change resamples a new target within the same cluster to simulate a change in user intent. Through this process, we enrich the benchmark with diverse forms of cross-turn dependency and interaction dynamics that go beyond simple sequential editing.

% \paragraph{Stage 3: Data Cleaning and Review.} 
% To ensure dataset quality, diversity, and consistency, we apply a combination of automatic filtering, balanced sampling, and manual review. First, we remove interactions with cyclic or repeated video transitions, such as cases where the retrieval sequence returns to an earlier video (e.g., $A \rightarrow B \rightarrow A$). We further filter out dialogues with semantically duplicate edits across turns by checking them with an LLM, ensuring that each turn contributes a meaningful update to the retrieval intent. To ensure broad semantic coverage, we categorize each sample along two dimensions: content category and modification type. Based on these labels, we perform approximately balanced sampling so that the benchmark spans diverse content domains and edit intents. 

% Finally, to ensure benchmark quality and verify that each session conforms to its predefined interaction scenario, we adopt a two-stage review protocol. Each sampled data entry is first screened by GPT-4o for automatic quality assessment. The retained samples are then evaluated by three independent human reviewers along three dimensions: accuracy, history dependency, and scenario consistency. To support consistent and standardized evaluation, we provide the reviewers with a detailed scoring guideline in App.~\ref{app:Scoring_Guideline}.

% This guideline ensures that all reviewers follow a consistent and standardized protocol when evaluating each sample, thereby improving the consistency and overall quality of the dataset. To enforce strict quality control, we retain a sample only when all three evaluation dimensions receive scores above 80. If exactly one dimension is scored below 80, the reviewers discuss the case and revise the user feedback accordingly, after which the sample is re-evaluated. Through this rigorous review process, the final benchmark achieves an average score of 85, further ensuring the quality of \Mybenchmark.

\subsection{Benchmark Construction} 
To instantiate the design principles above, we construct \Mybenchmark in three stages: initial multi-turn interaction generation, scenario-guided interaction construction, and data cleaning with quality review. The first stage produces seed interactions that are semantically coherent and sufficiently challenging for multi-turn retrieval. The second stage transforms these seed interactions into the predefined scenarios described in Sec.~2.1, thereby enriching the benchmark with diverse forms of cross-turn dependency. The final stage ensures benchmark quality and scenario consistency through multi-stage review.

\paragraph{Stage 1: Initial Multi-turn Interaction Generation.}
We begin with the WebVid-2M dataset and first partition the videos into 600 semantic clusters based on caption embeddings. This restricts reference--target sampling to locally coherent semantic neighborhoods, making the resulting modifications more natural and compositionally meaningful. For each reference video, we then sample a target from the same cluster under a similarity-range constraint, excluding candidates that are either too similar or too dissimilar. This encourages reference--target pairs that are suitable for progressive multi-turn refinement rather than trivial matching or abrupt semantic jumps. Given each reference--target pair, we prompt an MLLM to generate a user-style editing instruction that describes one or two salient modifications, leaving room for subsequent turns. To further control difficulty, we perform retrieval within the same cluster using a pretrained \task model. If the target can already be retrieved in a single turn, the pair is discarded as overly easy. Otherwise, the top-1 retrieved result is used as the next-turn reference, and the process is repeated by conditioning the MLLM on the accumulated dialogue history and the fixed target. In this way, we obtain seed multi-turn interactions that progressively approach the target video.

\paragraph{Stage 2: Scenario-Guided Interaction Construction.}
While the seed interactions generated in Stage 1 naturally capture progressive refinement, real-world interactive retrieval often exhibits more diverse forms of intent evolution. To instantiate the representative scenarios defined in Sec.~2.1, we further use GPT-4o to rewrite the seed interactions into scenario-specific dialogue trajectories. Specifically, \textbf{Rollback} is constructed by selecting a historical turn and rewriting the final feedback based on the reference video at that turn and the target video, so that the interaction explicitly returns to an earlier state (e.g., ``Compared to the most recent turn, I still prefer the video from turn \(x\). Building on that, ...''). \textbf{Combination} is formed by comparing both the current reference and a selected historical reference against the target video, and then composing a new feedback that jointly preserves constraints from both turns (e.g., ``Keep the ... from this turn and the ... from turn \(x\).''). \textbf{Correction} is created by rewriting the feedback in a self-corrective manner based on the current reference and the target video, using expressions such as ``Not ... but ...'', or ``Rather than ..., make it ...'' to locally rectify the retrieval trajectory. For \textbf{Preference Change}, we resample a new target video from the same cluster and generate a new feedback based on the current reference and this new target, thereby simulating a shift in user preference. Through this process, \Mybenchmark moves beyond simple progressive refinement and covers diverse forms of cross-turn dependency and intent evolution.

\paragraph{Stage 3: Data Cleaning and Quality Review.}
To ensure benchmark quality and scenario consistency, we perform data cleaning and multi-stage review. We first remove structurally invalid interactions, including trajectories with cyclic or repeated video transitions (e.g., \(A \rightarrow B \rightarrow A\)). We further discard dialogues in which any two feedback instructions are semantically equivalent, as determined by an LLM, ensuring that each turn contributes a meaningful update to the retrieval intent. To promote broad semantic coverage, we categorize each sample along two dimensions: content category and modification type. Based on these labels, we perform approximately balanced sampling so that the benchmark spans diverse visual domains and edit patterns, rather than being dominated by a narrow set of homogeneous scenarios. Finally, we adopt a two-stage review protocol to verify that each session conforms to its predefined interaction scenario and satisfies the desired quality criteria. Each sampled data entry is first screened by GPT-4o for automatic quality assessment. The retained samples are then evaluated by three independent human reviewers along three dimensions: accuracy, history dependency, and scenario consistency. To support consistent and standardized evaluation, we provide the reviewers with a detailed scoring guideline in App.~\ref{app:Scoring_Guideline}. We retain a sample only when all three evaluation dimensions receive scores above 80; otherwise, the reviewers revise the feedback and re-evaluate the sample. Through this rigorous process, the final benchmark achieves an average score of 85.


\paragraph{Statistic.}



\subsection{Evaluation Pipeline & Metric} 
Conventional retrieval metrics alone are insufficient for evaluating multi-turn composed video retrieval, as they only measure whether the target video is retrieved, but do not reveal whether the model has correctly accounted for the dialogue history when interpreting the final user intent. Therefore, in addition to standard retrieval metrics such as Recall@K and MRR, we further evaluate baseline methods from three complementary levels: Description-level Evaluation, which assesses the quality of the generated target description; Reasoning-level Evaluation, which measures the validity of the generated reasoning schema; and Retrieval-level Intent Satisfaction, which examines whether the top-k retrieved results satisfy the final retrieval intent.

For each evaluation level, model outputs are scored along three shared dimensions: \textbf{Intent Consistency}, \textbf{History Grounding}, and \textbf{Scenario Faithfulness}. Each dimension is rated on a 0--10 scale, where higher scores indicate better alignment with the final user intent, more appropriate use of the dialogue history, and stronger consistency with the predefined interaction scenario. To emphasize that a high-quality output must satisfy all three criteria simultaneously, we combine the three dimension scores multiplicatively within each evaluation level. Formally, for each level $l \in \{\mathrm{desc}, \mathrm{reason}, \mathrm{retr}\}$, the level-wise score is defined as
\[
S^{(l)} = s_{\mathrm{intent}}^{(l)} \cdot s_{\mathrm{hist}}^{(l)} \cdot s_{\mathrm{scen}}^{(l)},
\]
where $s_{\mathrm{intent}}^{(l)}$, $s_{\mathrm{hist}}^{(l)}$, and $s_{\mathrm{scen}}^{(l)}$ denote the 0--10 scores for Intent Consistency, History Grounding, and Scenario Faithfulness, respectively. The final score is then computed by summing the three level-wise scores across description, reasoning, and retrieval:
\[
S = \sum_{l \in \{\mathrm{desc}, \mathrm{reason}, \mathrm{retr}\}} S^{(l)}.
\]

To enable scalable fine-grained evaluation, we adopt an MLLM-as-a-Judge protocol in place of fully manual assessment. For each sample, the judge is provided with the full interaction history, the predefined scenario label, and the model output under evaluation, and assigns scores based on the criteria above. In this way, the protocol evaluates not only retrieval success, but also whether the model correctly reasons about multi-turn user intent with appropriate use of history and consistency with the underlying interaction scenario.